
So far we have used the computer to
drive the \robot\ in {\em interactive}\index{Interactive mode} or
{\em conversational}\index{Conversational mode} mode. In this mode of operation
we press a key on the computer keyboard, which passes
a command to the \robot, which carries it out immediately, unless by mistake
we try to make it leave the frame, in which case it answers with
a ``beep-beep''.
We go on like this, correcting driving mistakes as
the drawing takes shape, until we
tell the computer that the drawing is finished.

\smallskip

Whereas the interactive control of the \robot\ corresponds to the immediate
execution of a task, the programmed control of the \robot\ consists
in separating the description of the task from its later executions.

\smallskip

More
generally, the idea of {\em programming}\index{Programming} is the following:
to get something done, the {\em programmer} writes a {\em program}\index{Program}
(in a certain {\em language}\index{Programming language}) which describes precisely how
to do this thing. Then he asks a {\em computer}\index{Computer} to run
the program, following exactly the instructions given by the
program. In other words, instead of doing the thing
himself right away, the programmer explains to a computer how
to do it in his place later! In this way we avoid having to carry out
long and repetitive tasks, provided we know how to
describe these tasks properly.

\smallskip

If he has made a mistake, the programmer
must correct his errors by modifying the program. This is called
{\em debugging}\index{Debug} the program. The {\em debugging}\index{Debugging}
of the program consists in finding and then removing the {\em bugs}\index{Bug}, that is, the
errors which show up during execution as malfunctions (such as, for
example, {\idt qtl()} instead of {\idt qtr()}, or {\idt st()} forgotten at the end of the program)%
\footnote{\ In the French edition we use the vocabulary recommended by the French ministerial
terminology commission \cite{terminologie-info}, which translates {\em bug} by {\em bogue},
a French word that sounds the same. The American word {\em bug}
means small harmful insects and, more familiarly, germs
and other viruses, as well as the little things that go wrong, such as errors
in programs. Since the French {\em bogue} means the prickly husk
of the chestnut, the idea of nuisance is kept. The husk also
calls to mind a typographical slip. The French translation
does not keep the American word's idea of a little grain of sand with dramatic
consequences.}.
The principle of computers controlled by a program has
analogies with many human activities,
as the table \hbox{below shows:}

\begin{table}[htbp]
\begin{center}
\small
\begin{tabular}{@{}|p{25mm}|p{30mm}|p{30mm}|p{30mm}|@{}}
\hline
{\em Programmer } & {\em Program } & {\em Computer } & {\em Drawing }  \\
\hline \hline
Head chef & Recipe & Cooks & Meal\\
Com\-pos\-er & Musical work & Mu\-si\-cians & Con\-cert \\
Au\-thor & Text of a play  & Ac\-tors & The\-a\-tre per\-for\-mance \\
Cho\-re\-og\-ra\-pher & Cho\-re\-og\-ra\-phy & Dan\-cers & Bal\-let \\
Di\-rec\-tor & Film & Film pro\-jec\-tor & Film show\-ing \\
Wea\-ver & Punched cards & Jac\-quard loom & Jac\-quard fab\-ric\\
%Ma\^\i\-tre-Êcu\-yer & Tex\-te de re\-pri\-se & Ca\-va\-lier et che\-val & Re\-pri\-se \hfill de dres\-sa\-ge \\
Course de\-sign\-er & Course plan & Rider and horse & Show-jump\-ing round \\
\hline
\end{tabular}
\end{center}
%\label{tab:act-prog}
\caption{Examples of programmed activities}
\end{table}

We are now going to learn to program computers
in \python. Driving the \robot\ to make drawings will be our
main application. To make our drawings in
this programmed mode of operation, we proceed as follows:
\begin{enumerate}
\item We begin by writing a \python\ program in which we
state the sequence of commands that the \robot\ must
carry out to make the drawing we have in mind, or that
we have drawn beforehand on a sheet of squared
paper. This program is written with the help of the computer
(with a text editor), which keeps it on disk in a file whose name ends
with {\idt .py}, for example {\idt Square.py}.
\item  Next we ask a program of the computer, the \python\ {\em in\-ter\-pre\-ter}\index{Interpreter},
to read this program, to check that it is correctly written and to
translate it into the {\em machine language}\index{Machine language} of the computer
(or into a form very close to it). We say that the interpreter
{\em interprets}\index{Interpret} the program. The translated program and the original
\python\ program are equivalent in the sense that they describe
exactly the same sequence of commands for the \robot. They are
simply written in different languages.
\item  Finally the interpreter has the translated program {\em executed}\index{Execution}
by the computer. In doing so, the computer drives the \robot, asking it
to carry out the commands in the order in which we
wrote them in the original \python\ program.
\end{enumerate}
The last two steps are done with a single command, typed in
a terminal window (see the next chapter):
{\idt ./robot.sh Square.py} has the program {\idt Square.py} interpreted,
that is, checked, translated and then executed. You can also drag
the icon of the file {\idt Square.py} onto that of the application {\idt LeRobotPythonEN.app}.

{ %1
\newcommand{\touche}[1]{{
\begin{picture}(12,12)(0, 2)
\put(0,0){\framebox(12,12)[lb]{}}
\put(1.5,2){\makebox(0,0)[lb]{\tiny #1}}
\end{picture}\hspace{0.3ex}}}
\begin{exemple}[Square]{\rm
Suppose we want to draw the square opposite. The
\python\ program will be the following:

{\raggedright
 \vspace{\baselineskip}
\vspace*{1.3cm}%
\includevspacePostScript[-\baselineskip]{PetitCarre}{0.95}{0.95}%
\vspace*{-1.3cm}%
 \input python/Square.tex
 }
\noindent When it runs this program,
the computer passes to the \robot\ the order to carry out the sequence
 of commands {\idt fd()} then {\idt fd()} then
{\idt qtr()} then {\idt fd()} \ldots\ and finally {\idt qtr()}, exactly as if
we had typed on the keyboard \touche{F} then, one after the other, \touche{F},
\touche{+}, \touche{F}, \touche{F}, \touche{+}, \touche{F}, \touche{F},
 \touche{+}, \touche{F}, \touche{F} and finally \hbox{\touche{+}\@.} The last command
 {\idt st()} (from {\em stop}) tells the computer that it must wait until we press
 the escape key \touche{Esc} to end the execution of the
 program (another key makes it draw the drawing again). This leaves us time to admire the \robot's drawing
 at our leisure before it disappears. The drawing is not lost since
 we just have to run the program again to reproduce it.
 \noindent $\Box$ }\end{exemple}
 } %1

We can see from this first example that the \python\ language, although
designed in the Netherlands by Guido van Rossum at the {\sc Cwi} (the Centre for Mathematics
and Computer Science in Amsterdam), uses English
words such as `{\mcl from}' and `{\mcl import}':
the first line of the program imports from the library
{\idt LibRobot} all its commands, which is what the star `{\idt *}' stands for.
Computer science keeps the traces of its English-speaking
\hbox{origins!}

For the moment it is enough for us to know
that the computer that runs this program asks the \robot\ to carry out
the commands that come after the line `{\mcl from}~{\idt LibRobot}~{\mcl import}~{\idt *}',
one after the other, line after line, and from left to right on each line.
Each command is followed by {\idt ()}, which means that it needs
no further detail to be carried out. The
semicolons that separate the commands on the same line can be understood as
meaning `then'. There is no semicolon at the end of a line: going
on to the next line is enough to say `then'. The line that begins with the sign
`{\cmt \#}' is a comment, which the computer ignores (we shall come back to this in the next chapter).


%%%%%%%%%%%%%%%%%%%%%%%%%%%%%%%%%%%%%%%%%%%%%%%%%%%%%%%%%%%%%%%%%%%%%%%%%%%
\needspace{6\baselineskip}\noindent \rule{137mm}{0.1mm}\par\nopagebreak %---------------------------------------------
\begin{exercice}{
%\label{exo:trouver-cocotte}
Run the program {\idt Draw} for the interactive control
of the \robot\ and give it the orders described by
the following \python\ program (the correspondence between the keys of the
keyboard and the commands of the \robot\ is given in table
\ref{tab:cmd-inter} page \pageref{tab:cmd-inter}):
\input python/Unknown.tex
\noindent What drawing did you ask the \robot\ to make?
\noindent $\Box$}\end{exercice}
\begin{exercice}{
%\label{exo:corriger-E}
We would like the program below to draw the capital letter {\tt E}
\hbox{opposite:}

{\raggedright
\vspace{\baselineskip}
\vspace*{3cm}%
\includevspacePostScript[-\baselineskip]{E-epaisse}{0.95}{0.95}
\vspace*{-3cm}%
 \label{prog:LettreEavecBogues}
 \input python/LetterEWithBugs.tex
 }
\noindent This faulty program contains two bugs. Can you find them?
\noindent $\Box$}\end{exercice}
\begin{exercice}{
%\label{exo:exec-stegosaure}
Run the \python\ program {\idt Stegosaurus} found in the folder
{\idt book} that comes with this book.
\noindent $\Box$}\end{exercice}
%%%%%%%%%%%%%%%%%%%%%%%%%%%%%%%%%%%%%%%%%%%%%%%%%%%%%%%%%%%%%%%%%%%%%%%%%%%
\begin{corrige}{
%\label{cor:trouver-cocotte}
A paper bird (an origami ``cocotte'').
\noindent $\Box$}\end{corrige}
\begin{corrige}{
%\label{cor:corriger-E}
The program {\idt LetterEWithBugs} on page \pageref{prog:LettreEavecBogues}
contains two bugs:\\
\begin{minipage}[t]{115mm}
\begin{enumerate}
\item the command {\idt qtr()} on the third line
should be a command {\idt qtl()},
\item a command {\idt fd()} is missing on the last line but one.
\end{enumerate}
The result of running this program is shown opposite
(with the \robot\ in its final position).
The debugged program, which correctly draws the thick letter {\tt E}, is the following (the corrections are
un\-der\-lined):
\end{minipage}
\hfill
\begin{minipage}[t]{17mm}
\vspace{3.2cm}\includePostScript{lettre-e-avec-bogues}{0.95}{0.95}%
\hspace{-1.7cm}  % 1.7cm de large
\end{minipage}
\input python/LetterE.tex
\noindent $\Box$}\end{corrige}
\begin{corrige}{
%\label{cor:exec-stegosaure}
The \python\ program {\idt Ste\-go\-sau\-rus} is stored in the
file {\idt Stegosaurus.py} of the folder {\idt book/\allowbreak 04\_writing\_running}.
In a terminal window, opened in the folder {\idt LeRobot\_Python\_EN}, just
type the command
\begin{center}
{\idt ./robot.sh book/04\_writing\_running/Stegosaurus.py}
\end{center}
which has the \python\ program {\idt Stegosaurus.py} interpreted: the window of the \robot\ opens
and the stegosaurus is drawn. You get the same result by dragging the icon of the
file {\idt Stegosaurus.py} onto that of the application {\idt LeRobotPythonEN.app}.

\noindent $\Box$}\end{corrige}

%%%%%%%%%%%%%%%%%%%%%%%%%%%%%%%%%%%%%%%%%%%%%%%%%%%%%%%%%%%%%%%%%%%%%%%%%%%
\noindent \framebox[\textwidth]{
\begin{minipage}{\myparboxwidth}
\setlength{\parindent}{2em}
\small \sf
\vspace{\abovedisplayskip}
\section*{Addenda}
We owe to Seymour Papert, the author of {\sc Logo} \cite{logo}, the spreading of the idea
of using drawing to learn the basic concepts of programming. However,
the introduction of a grid spares us the use of angles and of lengths counted in
pixels\@%
\footnote{\ The {\em pixel}\index{Pixel} is the smallest
element of the graphics screen that can be lit (white), unlit (black) or
coloured, with more or fewer brightnesses and colours depending on
the quality of the screen.}. The grid makes it very easy to draw
arcs of
circles, which is not straightforward in {\sc Logo}.

The idea of using a
grid to make geometric drawings is not new. It has long been known
at school (see \cite{thevenet1} for example). It is found,
as the basis of graphics programming exercises, in
many high-level computing works, as in \TeX\@\footnote{\ \TeX, or rather its variant \LaTeX\
\cite{lamport1}, is the programming language we have used to
typeset this book.}\ for example (see \cite{knuth1} pages 389--390).
The grid
removes the problems of the discrete representation
of continuous mathematical concepts on a finite screen. As a result
our segments look like segments as we draw them on paper
(and not like lines put side by side in a staircase), and our
circles look like circles (and not
like squares, as happens on graphics screens for circles that are too
small).  Admittedly, the use of large squares limits the possibilities
of drawing. Programming is all the easier if we take advantage
of geometric shapes to guide the structure of programs.
Once the basic notions have been learnt, we shall be able to vary the size of the grid
so that, going to the limit set by the pixels of the screen,
we recover all the possibilities offered by today's high-resolution
graphics screens.

In the next chapter we begin to write our first
\python\ programs. The algorithms\@\footnote{\ An {\em algorithm} is the
description of a sequence of operations on objects to
accomplish a task. This word comes from the Latin version of the name of the
Persian mathematician and astronomer al-Khwarizmi of the 9th
century AD.
His treatise {\em Al-Jabr wa l-muqabala} teaches how to solve algorithmically
equations of the first and second degree with numerical coefficients.
The term {\em algebra\/} derives from the very title of this treatise.
} remain very simple since they
consist in codifying the movement of the hand that draws.

\vspace{\abovedisplayskip}
\end{minipage}
}
%%%%%%%%%%%%%%%%%%%%%%%%%%%%%%%%%%%%%%%%%%%%%%%%%%%%%%%%%%%%%%%%%%%%%%%%%%%
